\section{Introduction}
The topic of gesture-based gaming on a touch-device spans a number of distinct research areas. This chapter briefly covers the most relevant of these fields. This includes multi-touch interaction where the various ways of controlling a touch-device are discussed. The use of gestures, in contrast to other more conventional means of input, is also considered. As these techniques will be applied to gaming, a discussion of the ways in which gaming experience can be evaluated is also included. Finally, the limited research on gesture-based gaming is reviewed.

\section{Multi-touch}
\subsection{Introduction}
The way in which we interact with computers has changed very little over the last few decades. That is, the keyboard and mouse remain the primary means through which individuals interact with graphical user interfaces~\citep{Kin09}. However, these means of interaction greatly limit the number of simultaneous input points and consequently the types of interfaces that can be developed~\citep{Mosc08}. In addition, these devices distance the user by limiting the degree and magnitude with which they can physically interact with the system~\citep{Mosc06}. Multi-touch technology, however, offers an alternative to the traditional computing paradigm. Multi-touch, in relation to touch-screen devices, refers to the ability to simultaneously control two or more input points or touches~\citep{Jaim07}. These systems allow users to physically interact with the computing device with the option to use two or more input points. This in turn, increases the degrees of freedom that the user is able to manipulate, which has implications for the type of applications that are possible and the way in which interaction itself is performed~\citep{Mosc08}.


\subsection{Types of interaction}
At present, there is no standard hardware with which \gls{c2:multitouch} interfaces are developed. Some devices utilise capacitive technology, while others use cameras and Computer Vision techniques in order to provide touch capabilities~\citep{Wigd09}. This results in there being a number of variations in the input techniques in use. These techniques include direct-touch versus indirect-touch, single-hand versus bimanual, as well as the use of differing numbers of fingers~\citep{Kin09}. There are also techniques where the shape and orientation of the user's whole hand is tracked and identified~\citep{Wu03}. These mechanisms will be discussed in further detail below.

\subsubsection{Direct-touch versus indirect-touch}
Indirect-touch refers to the use of a pointing device to control and manipulate objects onscreen~\citep{Mosc08}. Users are required to visually track this pointer or cursor. Direct touch, in contrast, refers to directly touching or interacting with the object one wishes to manipulate~\citep{Kin09}. Thus one tracks the finger or fingers interacting with the object rather than an onscreen cursor~\citep{MoscH06}. 

Perhaps the most well known example of \gls{c2:indirect} is a mouse. However, research has also been conducted on using multiple fingers in order to control a number of control points in an indirect manner~\citep{MoscH06,Mosc08}. In this scenario one manipulates multiple onscreen control points through the use of more than one finger. However, this manipulation is done on a device connected to a separate output screen rather than on the display screen itself~\citep{MoscH06}.

Direct-touch has been shown to be faster than using single-point, indirect-touch methods, such as a mouse~\citep{Kin09}. However, no comparison for speed has yet been performed between direct-touch and indirect, multi-touch methods. Direct-touch, however, suffers from issues of occlusion. That is, when interacting with an object, one often hides or blocks portions of the object~\citep{Kin09}. This is of concern for interfaces with small buttons and other such elements. In addition, due to the size of a finger relative to that of a mouse cursor, both direct and indirect touch-based methods suffer from decreased precision of selection~\citep{Kin09}. However, as mentioned above, the nature of touch-based interaction allows for a greater number of control points, and thus degrees of freedom within an interaction, when compared to conventional mouse-based pointing devices.

\subsubsection{Multi-finger techniques}
Many touch-devices support multiple points of input. This increases the degrees of freedom that a user has when manipulating or controlling an object or interface~\citep{Mosc08}.

However, research has found that using multiple fingers on the same hand often produces a minimal increase in the degrees of freedom the user is able to control~\citep{Mosc08}. This is due to the fact that fingers on the same hand are linked physically and neurologically~\citep{Mosc06}. Thus the movements produced by multiple fingers are often highly similar. Therefore, while multiple fingers on the same hand may increase input bandwidth on tasks that require coordinated control, such as using two fingers to scale an object, they do not actually offer a large increase in the degrees of freedom available to a user~\citep{Kin09}. They may, in fact, hinder a task as they increase the chance of simultaneously selecting multiple targets. Finally, research has suggested that users often do not attempt to use more than two fingers simultaneously unless prompted to do so~\citep{Kin09}. This suggests that while many devices offer recognition of multiple simultaneous input points, this multiplicity is often not used or required in unimanual applications.

\subsubsection{Unimanual versus bimanual interaction}
Devices that are able to detect multiple simultaneous inputs also present the opportunity for interaction with more than one hand. This too can increase the degrees of freedom available to a user~\citep{Kin09}.

Bimanual interaction does not suffer from the same limitations as unimanual, multi-finger interaction due to the decoupled nature of the input points~\citep{Mosc08}. That is, as the multiple input points come from two hands, their motion is far more separable. However, this decoupling of movement also means that it is more difficult to perform coordinated tasks using two hands instead of one~\citep{Mosc08}. Bimanual interactions have also been shown to have a higher miss rate than unimanual conditions~\citep{Kin09}. \citet{Mosc08} showed that bimanual interaction is preferable when performing separable tasks but that unimanual (multi-finger) interaction is preferred when attempting to perform more precise, coordinated tasks. A further constraint on bimanual manipulation is that, without a clear correspondence between each hand and the control points, users often become confused. This may be due to the increased cognitive load of having to control and coordinate both hands simultaneously~\citep{Kin09}. With this said, however, bimanual interaction shows a speed advantage over unimanual interaction in decoupled tasks that do not require coordinated movement~\citep{Kin09}.

The differences in the degrees of freedom, accuracy and control presented by each method suggests that these variants are not interchangeable. Thus it appears that while both unimanual and bimanual interaction have their benefits, the type of application and level of control required dictates which of the two methods is more suitable.

\subsection{Benefits}
As is apparent, there are a number of ways in which touch-devices can be utilised. Furthermore, there does not appear to be one universally ideal technique. Instead it seems that the type of interaction style that should be used is highly application dependant. With this said, multi-touch technology does offer a number of distinct benefits over traditional interaction means such as the mouse and keyboard.

One such advantage is that, in many cases, applications that utilise touch-devices are often easier for uninitiated users to learn. As shown by \citet{Mosc06}, all that is usually needed is the instruction to touch the device. This in turn makes applications that use touch interaction far more accessible. This may be because multi-touch interfaces have been shown to be more naturalistic, which perhaps aids learning of a novel system~\citep{Forl07}.

With regard to speed, it has been shown that the fastest multi-touch interaction (bimanual, direct-touch) is roughly twice as fast as using a mouse pointer. Furthermore, the use of multiple fingers or hands enables task parallelism and thus increases task performance~\citep{Kin09}. The increase in the number of input points and degrees of freedom also means that input points can be used simultaneously to control more than just current position. One example is the use of multi-touch to combine translation and rotation of an object into one operation by utilising multiple fingers~\citep{Mosc06}.

Perhaps the most important benefit of multi-touch devices, and the interaction techniques they make possible, is that they allow for a whole range of novel applications and ways in which one can interact with a computing device.

\subsection{Limitations}
As with any interaction technique, multi-touch has certain limitations when compared to other input mechanisms. These limitations are primarily caused by the nature of touch computing, that is, users use their fingers, rather than a more precise pointer (such as a mouse cursor), to interact with objects on screen~\citep{Mosc06}.

Direct-touch suffers from issues of occlusion not present in indirect methods~\citep{Wu03}. This has implications for the design of applications that utilise touch technology. While some work has been done on mediating these issues by making use of techniques such as using a second finger to manipulate the scale or coordinate system of the display, occlusion remains a problem within touch interface design~\citep{Mosc08}.

Another issue related to the size of a finger when compared to other pointing devices, is that of accuracy. Touch-based interaction has been shown to have higher miss rates when compared to mouse-based pointers~\citep{Mosc08}. In addition, the precision of movements available is greatly reduced when compared to a mouse or pen device. This reduced accuracy is also not fully eliminated by the use of indirect-touch methods~\citep{MoscH06}. Furthermore, due to the nature of interaction, there are limitations on the proximity and size of the display~\citep{Mosc08}. The display also fails to provide any form of haptic feedback. For example, touch typing would be considerably more difficult, as one is unable to feel the relation of different keys to each other. However, this is often also lacking in other interaction techniques~\citep{Wobb09}.

Finally, the use of more than one pointing device can also, in some cases, cause the user to become disoriented and confused due to the increased cognitive load~\citep{Mosc08}. The same has been shown to be true for multi-finger interaction when this movement is not coordinated~\citep{Kin09}.

\subsection{Conclusion}
Multi-touch technology presents the opportunity for creating applications that utilise two or more control points. This has been shown to increase task parallelism and selection speed as well as changing the way in which people are able to manipulate and control digital objects and interfaces. However, despite these benefits, touch technology also introduces issues of occlusion and inaccuracy. With regard to cognitive performance, while touch technology has been shown to be highly intuitive to the uninitiated, problems of increased cognitive load are introduced when using multiple fingers or hands. Therefore, it appears that for multi-touch technology to be truly effective, the application itself must suit this style of interaction. Thus touch interaction should not be viewed as a panacea. Instead one should see it as a technique that, within the right framework, can greatly improve the quality of interaction of an application.


